% =====================================================================
%  Related work (ban nhap Claude 20-11; chu nhan / thay Van sua giong van).
%  Nguon: paper/paper1/sections/02_background.tex (~1750 tu). Muc tieu ~900 tu.
%  Da ap: 1.2 (bo "concentration ... destroys it", khong nhan qua); 1.6 (Cirillo:
%  chi giu dieu da doc trong toan van -- bang 1.6 cua docs/PAPER1_GD1_KET_QUA.md;
%  bo "suy giam sau ~60 CNOT", "ho quy cho loi cong", "94.60 / 94.13 %" vi chua
%  doi chieu lai); R2 (dong gop la giao thuc benchmark, khong nhan loi the).
%  Bo han muc "Bounded advantage windows" (Carducci 2026): nhom chi doc duoc tom
%  tat, luat chu nhan: tom tat = chua kiem. Muc carducci2026 trong bibliography
%  con lai -> go khi soat danh muc tham khao.
%  Bo doan "We take the reviewers' point ..." (AE che nhac phan bien).
%  Gillani 2026: giu su that cua ban TETC (doc toan van luc do), rut gon; chua
%  doi chieu lai trong GD 2.
% =====================================================================
\section{Related work}\label{sec:related}

\subsection{Quantum kernels}

A quantum kernel maps $\bm{x}$ to a state $|\phi(\bm{x})\rangle$ prepared by a parameterised
circuit and uses the fidelity between states as a similarity \cite{havlicek2019,schuld2019}.
The construction belongs to a broader programme of quantum machine learning
\cite{ref5,ref6,ref7}, and the correspondence between parameterised-circuit models and kernel
machines \cite{ref9} lets the quantum part be placed inside a classical learner. Separations
from classical learners have been proved for constructed problems \cite{ref18}, but the
conditions under which they carry over to data from practice are restrictive \cite{ref19,ref20}.

Two lines of theory bear on what to expect. Alignment-based bounds \cite{cristianini2002,cortes2012}
state that, at equal empirical risk, a kernel better aligned with the labels admits a tighter
generalisation guarantee; they do not say whether a given circuit produces such a kernel.
Thanasilp \emph{et al.} \cite{thanasilp2024} prove that fidelity kernel values concentrate
exponentially in the number of qubits under conditions that include high expressivity,
entanglement and noise. The effect is related to the barren plateaus of variational training
\cite{ref10} and to expressivity \cite{ref21}: the properties that make a feature map
expressive can also make its Gram matrix uninformative at large width. We measure alignment
and Gram-matrix dispersion across widths (Sec.~\ref{sec:res-width}) and report how they vary
with macro-$F_1$; we do not test the conditions of that theorem.

\subsection{Quantum machine learning for intrusion detection}

Quantum kernels and variational classifiers have been applied to network intrusion detection
in a growing number of studies \cite{ref13,ref14,ref15}, most on NSL-KDD \cite{ref4} or a later
corpus \cite{ref25}. A 2026 survey \cite{kaissar2026} finds that most approaches rely on
simulated quantum environments and older datasets and that evaluation practice is inconsistent
across studies. A common design evaluates one qubit count, one training-set size and one class
prior, and compares accuracy with a small set of classical baselines.

The closest study is the benchmark of Cirillo \emph{et al.} \cite{qmi2026}. Its main dataset is
ToN\_IoT, and the best configurations are re-run on NSL-KDD. It evaluates Pegasos-QSVC over
$378$ configurations and $14$ feature maps, a variational classifier and a hybrid
quantum--classical network, in ideal simulation and under noise models derived from three IBM
backends with $1024$ shots, and compares them with an SVM, XGBoost, a random forest and a neural
network at matched dimension. Results are means and standard deviations over four seeds; the
study states that it does not aim to demonstrate quantum advantage. Several elements of our
design were already present there: a quantum SVM for intrusion detection on NSL-KDD with
simulated backend noise, the comparison of a fully entangled \ZZ{} map with the
non-entangling \Zmap{} map, projection to four qubits, and tree-ensemble baselines on the
same number of dimensions. We do not claim these as new. Two of its findings agree with ours.
The classical models keep the higher peak score, and the best balance between performance and
cost lies at $20$ to $35$ CNOT gates; our four-qubit circuit uses $24$.

Our design differs from that study in five ways. It sweeps the training-set size from $10^2$
to $10^4$, where that study uses one class-balanced subsample ($15{,}114$ training records on
NSL-KDD). It tests every comparison with paired runs and multiplicity correction, where that
study reports means and standard deviations without a test. It tunes quantum and classical
models by the same procedure and separates three classical kernels. It measures how the Gram
matrix changes with circuit width. And it adds UNSW-NB15 and an analysis of test records that
duplicate training records.

\subsection{Benchmarks of quantum models}

Bowles, Ahmed and Schuld \cite{bowles2024} evaluate twelve quantum models on $160$ mostly
synthetic datasets in noise-free simulation with up to $18$ qubits, against a multilayer
perceptron, an RBF-kernel SVM and a convolutional network, and find that the quantum models
rarely outperform well-tuned classical baselines. Schnabel and Roth \cite{schnabel2024} train
more than $20{,}000$ tuned models comparing fidelity and projected quantum kernels on five
families of datasets with up to $15$ qubits, against RBF-kernel baselines, and use correlation
analysis to find which design choices matter. Neither varies the training-set size: Bowles
\emph{et al.} use fixed counts per task, and Schnabel and Roth use $240$ training and $60$
test points throughout.

Gillani \emph{et al.} \cite{gillani2026} benchmark quantum and classical models for intrusion
detection on four corpora, including NSL-KDD and UNSW-NB15, at eight qubits with an ablation
over $4$, $6$ and $12$, against five tuned classical baselines including a random forest and
XGBoost, with Benjamini--Hochberg control over $108$ comparisons on NSL-KDD. They conclude that
reported quantum gains are largely explained by the classical preprocessing once the same
preprocessing is given to the classical models. Our results agree that the front-end carries
most of the performance: given all $122$ features, XGBoost is not beaten by the four-qubit
kernel at any size we tested (Sec.~\ref{sec:res-n}). They differ on the four-component
embedding, where the kernel is significantly ahead of both tree ensembles at $N=10^4$. Their
comparison uses eight qubits, where our kernel loses every comparison (Sec.~\ref{sec:res-width}),
and the full training partition, which our sweep does not reach. The two designs therefore do
not overlap where our lead appears, and our data cannot say which difference accounts for the
gap.

\subsection{Positioning}

% Bang dinh vi (go tay). Lay tu paper/paper1/tables/novelty_matrix.tex (o nao cung
% doc tu toan van bai duoc trich, theo ban TETC); 20-11 sua cho QMI: caption bo
% "reviewers pointed us to" va Carducci (chi doc duoc tom tat); cot This work ghi
% nhieu "mot backend" va hieu chinh da so sanh moi (BH, Holm ca ban do).
% Chua doi chieu lai toan van Bowles / Schnabel / Gillani trong GD 2.
\begin{table*}[t]
\centering
\caption{Positioning against the closest benchmarking studies, read from their full texts;
``n/r'' means not reported. The last three rows are the axes on which this work differs; on
the rows above it is comparable or narrower.}
\label{tab:positioning}
\footnotesize
\setlength{\tabcolsep}{4pt}
\begin{tabular}{@{}l p{2.05cm} p{2.05cm} p{2.05cm} p{2.35cm} p{2.6cm}@{}}
\toprule
 & \textbf{Bowles et al.\ \cite{bowles2024}} & \textbf{Schnabel \& Roth \cite{schnabel2024}} & \textbf{Cirillo et al.\ \cite{qmi2026}} & \textbf{Gillani et al.\ \cite{gillani2026}} & \textbf{This work} \\
\midrule
Domain          & general QML        & general QKM        & network IDS        & network IDS        & network IDS \\
Datasets        & 6 tasks / 160 sets & 5 families / 64 sets & NSL-KDD, ToN\_IoT & NSL-KDD, UNSW-NB15, CICIDS2017, NF-ToN-IoT-v2 & NSL-KDD, UNSW-NB15 \\
Test split      & synthetic splits   & fixed 240/60       & random 80/20 on a 15\% class-balanced subsample & official splits    & official splits \\
Quantum model   & 12 QML models      & FQK, PQK           & Pegasos-QSVC, VQC, HQNN & hybrid VQC, IQP fidelity \& projected kernels & ZZ / Z fidelity kernels \\
Qubits          & up to 18           & up to 15 ($=$ features) & 2--10 swept, best 7 & 8 (abl.\ 4/6/12)   & 4 (NSL), 6 (UNSW), selected by rule \\
RF + XGBoost baselines & no          & no                 & yes                & yes                & yes \\
Repeats         & n/r                & $>$20{,}000 models (HPO) & 4 seeds      & 3--5 seeds         & 10 paired runs \\
Multiplicity control & n/r           & correlation analysis & none (mean $\pm$ sd) & BH over 108 comparisons & Holm within family; BH and Holm over the map \\
Noise model     & noise-free         & noise-free         & ideal + three noisy simulators & NISQ noise sweep & ideal; one noisy simulator \\
\midrule
\textbf{Training-set size swept} & no & no                & no                 & no (width ablation only) & \textbf{yes}, $N=10^2$--$10^4$ \\
\textbf{Verdicts by condition} & aggregate ranking & aggregate correlations & $F_1$ vs.\ CNOT count & per-dataset verdict & \textbf{109 paired verdicts} over prior, composition, shift, size \\
\textbf{Width rule transferred} & no & no & no            & no                 & \textbf{yes}: $n^\ast=4\to6$, thresholds unchanged \\
\bottomrule
\end{tabular}
\end{table*}


Table~\ref{tab:positioning} compares this work with the four studies above. On model, datasets,
circuit width and noise treatment, this work is comparable or narrower. It differs in three
respects: it sweeps the training-set size, it reports a paired verdict for each condition, and
its width rule is applied unchanged to a second dataset. The training-set size matters because
the ranking between \QSVM{} and the strong classical baselines reverses inside
$N\in[10^2,10^4]$ under the natural prior, so a benchmark at a single $N$ can report either
outcome. The contribution of this paper is that protocol and the map of verdicts it produces.
